% answers_short.tex —— 1998 年·上海卷·文史类 · 简短答案速查
% 由 tools/make_short_answers.py 自 answers.tex 的 \gksolhead/\gkanswer 提取，
% 并与 1998/1998_数学真题_上海文_答案要点.md 逐题比对后生成，勿手改（改源 MD 后重跑）。

\begin{gknotice}
\textbf{使用说明}：本节只给最终答案便于快速核对，完整推导见解析版。答案为固定来源 PDF 卷内「答案及评分标准」所载（出版者/版次未核实），非官方标准答案；第 15 题按原页核读为 $11\in B$。
\end{gknotice}

\gkpart{填空题}
\begin{itemize}[leftmargin=2.6em]
  \item[\textbf{1.}] $2$。
  \item[\textbf{2.}] $5$。
  \item[\textbf{3.}] $4$。
  \item[\textbf{4.}] $f^{-1}(x)=(x-2)^3+1$（$x\in\mathbb R$）。
  \item[\textbf{5.}] $\dfrac{2\sqrt2}{3}$。
  \item[\textbf{6.}] $11$ 天。
  \item[\textbf{7.}] $\dfrac{18}{35}$。
  \item[\textbf{8.}] $4$。
  \item[\textbf{9.}] $4$。
  \item[\textbf{10.}] $2$。
  \item[\textbf{11.}] $\dfrac{x^2}{9}-\dfrac{y^2}{16}=1$。
\end{itemize}

\gkpart{选择题}
\begin{center}\begin{tabular}{cccccc}
\toprule
题号 & 12 & 13 & 14 & 15 & 16 \\
\midrule
答案 & B & A & C & D & C \\
\bottomrule
\end{tabular}\end{center}

\gkpart{解答题}
\begin{itemize}[leftmargin=2.6em]
  \item[\textbf{17.}] $\sin\left(\dfrac{37\pi}{6}-2\alpha\right)=\dfrac{7+24\sqrt3}{50}$。
  \item[\textbf{18.}] $z_2=6+2i$，$|z_2|=2\sqrt{10}$。
  \item[\textbf{19.}] （1）二面角 $O_1\text{-}BC\text{-}D$ 的大小为 $\arctan(2\sqrt3)$；（2）异面直线 $OB_1$ 与 $D_1E$ 所成角为 $\arccos\dfrac{7\sqrt{10}}{40}$。
  \item[\textbf{20.}] （1）$y^2=4(x-1)$；（2）$\alpha=\dfrac\pi6$ 或 $\dfrac{5\pi}6$。
  \item[\textbf{21.}] $T(t)=t^3-3t+60$；11:00 与 14:00 均为最高温度 $62^\circ\mathrm C$；8:00–16:00 平均温度 $60^\circ\mathrm C$。
  \item[\textbf{22.}] （1）$b_n=-\dfrac{12n+5}{4}$；（2）$\lim\limits_{n\to\infty}\dfrac{k_1+k_2+\cdots+k_n}{a_nb_n}=\dfrac13$；（3）$c_n=-12n-89$。
\end{itemize}
